% Clase 9 --- El campo del dipolo obtenido DERIVANDO el potencial (§1.2).
% En Clase 4 el campo se calculó sólo sobre el plano bisector, sumando
% vectores. Acá se parte de V (que Clase 8 dio en un punto arbitrario) y se
% deriva: sale el campo en CUALQUIER punto, con mucho menos trabajo.
\begin{center}
\begin{tikzpicture}[scale=1]

  % =============== (a) camino largo: sumar vectores ===============
  \begin{scope}
    \node[figpos] at (0,0.55) {}; \node[figlbl, anchor=east] at (-0.12,0.55) {$+Q$};
    \node[figneg] at (0,-0.55) {}; \node[figlbl, anchor=east] at (-0.12,-0.55) {$-Q$};
    \node[figneutra, minimum size=5pt] (P) at (2.3,0.9) {};
    \node[figlbl, anchor=south] at (2.3,1.05) {$P$};
    \draw[figvec] (P) -- ++(15:1.0);
    \draw[figvec] (P) -- ++(-160:0.85);
    \node[figlblaux, anchor=west] at (3.3,1.15) {$\vec E_+$};
    \node[figlblaux, anchor=north east] at (1.5,0.42) {$\vec E_-$};
    \node[figlbl, align=center] at (1.4,-1.85)
      {sumar \textbf{vectores}\\ (y sólo sale fácil en el bisector)};
    \node[figlbl] at (1.4,-2.9) {(a) camino de la Clase 4};
  \end{scope}

  % =============== (b) camino corto: derivar V ===============
  \begin{scope}[xshift=7.4cm]
    \node[figpos] at (0,0.55) {}; \node[figlbl, anchor=east] at (-0.12,0.55) {$+Q$};
    \node[figneg] at (0,-0.55) {}; \node[figlbl, anchor=east] at (-0.12,-0.55) {$-Q$};
    \node[figneutra, minimum size=5pt] (P2) at (2.3,0.9) {};
    \node[figlbl, anchor=south] at (2.3,1.05) {$P$};
    \node[figlblaux, figblue, anchor=west] at (2.5,0.85) {$V(r,\theta)$};
    \draw[figvecres, line width=1.4pt] (P2) -- ++(-38:1.15);
    \node[figlbl, figred, anchor=west] at (3.3,0.15) {$\vec E = -\nabla V$};
    \node[figlbl, align=center] at (1.4,-1.85)
      {derivar un \textbf{escalar}\\ (vale en cualquier punto)};
    \node[figlbl] at (1.4,-2.9) {(b) camino de esta clase};
  \end{scope}

\end{tikzpicture}\\[0.5em]
{\small\itshape El contenido físico es el mismo, pero el orden de trabajo
cambia todo. Calcular $\vec E$ directo obliga a sumar vectores punto por punto;
calcular primero $V$ ---que es un número--- y después derivarlo da el campo en
\textbf{cualquier} punto de una sola vez. Las derivadas parciales que aparecen
($\partial V/\partial x$, $\partial V/\partial y$, \dots) son justamente las
componentes de $\vec E$ cambiadas de signo.}
\end{center}
